% 論文節の日本語訳（読み物 PDF）用 preamble
% lualatex + jlreq
\documentclass[12pt,paper=a4]{jlreq}
\usepackage{luatexja-fontspec}
\usepackage{amsmath,amssymb,bm}
\usepackage{graphicx}
\usepackage{hyperref}
\usepackage{enumitem}
\usepackage{booktabs}
\usepackage{microtype}
\usepackage{xcolor}
\usepackage{caption}
\usepackage{tcolorbox}
\tcbuselibrary{breakable,skins}

\hypersetup{
  colorlinks=true,
  linkcolor=blue!60!black,
  urlcolor=blue!60!black,
  citecolor=blue!60!black,
}

\graphicspath{{images/}{./}}

\setlength{\parindent}{1em}
\setlength{\parskip}{0.55em}
\setlist[itemize]{leftmargin=1.6em,itemsep=0.25em,topsep=0.35em}
\setlist[enumerate]{leftmargin=1.8em,itemsep=0.25em,topsep=0.35em}

\captionsetup{
  font=small,
  labelfont=bf,
  skip=0.4em,
  justification=raggedright,
  singlelinecheck=false,
}

\newcommand{\vx}{\bm{x}}
\newcommand{\vu}{\bm{u}}
\newcommand{\veps}{\bm{\epsilon}}
\newcommand{\vw}{\bm{w}}

\DeclareMathOperator*{\argmax}{arg\,max}
\DeclareMathOperator*{\argmin}{arg\,min}
\DeclareMathOperator{\Tr}{Tr}
\DeclareMathOperator{\diag}{diag}

% 命題・定理ボックス（ページ跨ぎ時も切れ目を抑える）
\newtcolorbox{propositionbox}[1][]{%
  enhanced,
  breakable,
  colback=blue!4!white,
  colframe=blue!45!black,
  coltitle=white,
  fonttitle=\bfseries\small,
  title={#1},
  attach boxed title to top left={yshift=-2mm, xshift=3mm},
  boxed title style={
    colback=blue!45!black,
    colframe=blue!45!black,
    boxrule=0pt,
    arc=1pt,
  },
  boxrule=0.5pt,
  arc=1.5pt,
  left=5pt, right=5pt, top=8pt, bottom=5pt,
  before skip=0.9em,
  after skip=0.9em,
  pad at break=3mm,
}

% 図はフロートせず本文位置に置く（分断を防ぐ）
\newcommand{\translationfigure}[2]{%
  \begin{center}
    \includegraphics[width=0.88\linewidth]{#1}\\[0.35em]
    {\small #2\par}
  \end{center}
  \vspace{0.4em}%
}

\begin{document}
{\color{gray}\footnotesize 出典: §1.3\par}
\vspace{0.2em}

\section*{1.3 モデリングの分類}


これまで見てきたように、DGM は広いスペクトルのモデリング戦略にまたがる。根本的な区別は、これらのモデルが基礎となるデータ分布をどうパラメータ化するか——すなわち、学習目的とは独立に、$p_\phi(x)$ を明示的に指定するか、それとも暗黙的にのみ指定するか——にある。

\begin{itemize}
\item \textbf{明示的モデル（Explicit Models）}: これらのモデルは、扱いやすい、あるいは近似的に扱いやすい密度／質量関数を通じて確率分布 $p_\phi(x)$ を直接パラメータ化する。例として AR、NF、VAE、DM があり、いずれも $p_\phi(x)$ を厳密に、あるいは扱いやすい下界を通じて定義する。
\end{itemize}


\begin{itemize}
\item \textbf{暗黙的モデル（Implicit Models）}: これらのモデルは、典型的にはあるノイズ変数 $z\sim p_{\mathrm{prior}}$ に対する $x=G_\phi(z)$ という形のサンプリング手続きを通じてのみ分布を指定する。この場合、$p_\phi(x)$ は閉形式では得られず、そもそも定義されないこともある。
\end{itemize}


表1.1は、これらの対照的なアプローチの簡潔な要約を与える。

\translationfigure{table-1-1.jpg}{表1.1: 明示的生成モデルと暗黙的生成モデルの比較}


\textbf{拡散モデルとの接続。} まとめると、これらの古典的な DGM の族は、複雑な分布をモデル化するための相補的な戦略を例示する。単独での重要性を超えて、それらは拡散モデルを理解するための指針も提供する。拡散手法はこれらの視点のいくつかから着想を受け継ぐ。すなわち、変分学習目的を通じて VAE と、対数密度の勾配（エネルギー関数と密接に結びつく）を学習するスコアマッチング手法を通じて EBM と、連続時間変換を通じて NF と接続する。

後の章で議論する拡散手法の基礎を据えるため、三つの中心的パラダイム——VAE（第2.1節）、EBM（第3.1節）、NF（第5.1節）——に焦点を当てる。この探求は、続く章でさらに展開される現代の拡散ベース生成モデリングの根底にあるコア原理への基礎を与える。
\end{document}
